\documentclass[12pt,a4paper]{article}

% ---- Encoding & Font ----
\usepackage{iftex}
\ifPDFTeX
  \usepackage[utf8]{inputenc}
  \usepackage[T1]{fontenc}
\else
  \usepackage{fontspec}  % XeTeX/LuaTeX (tectonic): Unicode nativo (·, —, ¿, ª, §)
\fi
\usepackage[spanish]{babel}
\usepackage{csquotes}

% ---- Page layout ----
\usepackage[top=2.5cm, bottom=2.5cm, left=2.5cm, right=2.5cm]{geometry}
\usepackage{setspace}
\onehalfspacing

% ---- Math ----
\usepackage{amsmath, amssymb, amsthm}
\usepackage{mathtools}
\usepackage{physics}
\usepackage{esint}  % \oiint

% ---- Graphics ----
\usepackage{graphicx}
\usepackage{tikz}
\usepackage{float}
\usepackage{caption}

% ---- Tables ----
\usepackage{array}
\usepackage{booktabs}
\usepackage{tabularx}
\usepackage{colortbl}

% ---- Colours ----
\usepackage{xcolor}
\definecolor{notebg}{HTML}{FFF3CD}
\definecolor{noteborder}{HTML}{FFC107}
\definecolor{boxred}{HTML}{F8D7DA}
\definecolor{boxredborder}{HTML}{F5C6CB}

% ---- Boxes ----
\usepackage{mdframed}
\usepackage{tcolorbox}
\tcbuselibrary{most}

\newtcolorbox{keybox}[1][]{
  colback=blue!5,
  colframe=blue!70!black,
  arc=4pt,
  boxrule=1pt,
  fonttitle=\bfseries,
  title={#1}
}

\newtcolorbox{warnbox}[1][]{
  colback=boxred,
  colframe=boxredborder,
  coltitle=black!75,
  arc=4pt,
  boxrule=1pt,
  fonttitle=\bfseries,
  title={#1}
}

\newtcolorbox{notebox}[1][]{
  colback=notebg,
  colframe=noteborder,
  coltitle=black!75,
  arc=4pt,
  boxrule=1pt,
  fonttitle=\bfseries,
  title={#1}
}

% ---- Hyperlinks & TOC ----
\usepackage[hidelinks]{hyperref}
\usepackage{bookmark}

% ---- Enumerate ----
\usepackage{enumitem}

% ---- Miscellaneous ----
\usepackage{icomma}
\usepackage{siunitx}
\usepackage{textcomp}
\usepackage[bottom]{footmisc}

% ---- Diagramas (gráficas y esquemas) ----
\usepackage{pgfplots}
\pgfplotsset{compat=1.18}
\usetikzlibrary{babel}  % babel español activa `>`; esto evita romper las flechas `->`
% courses/ElecMag2024/Clases/assets/tikzstyles.tex
% ---------------------------------------------------------------------------
% Estilos compartidos para los diagramas del curso Electromagnetismo (2024).
% Fuente única del "look" visual (colores, grosores, escalas) para que todas
% las figuras (TikZ / pgfplots / CircuiTikz) sean consistentes entre clases.
%
% Es una COPIA de courses/Fisica3-2015/Clases/assets/tikzstyles.tex: mismo
% dominio, misma paleta. Copia y no symlink para que cada curso sea
% autocontenido y la paleta de Física III pueda evolucionar aparte.
%
% Uso: la clase debe cargar antes `tikz` y `pgfplots` (y `circuitikz` si la
%      clase tiene circuitos), y luego % courses/ElecMag2024/Clases/assets/tikzstyles.tex
% ---------------------------------------------------------------------------
% Estilos compartidos para los diagramas del curso Electromagnetismo (2024).
% Fuente única del "look" visual (colores, grosores, escalas) para que todas
% las figuras (TikZ / pgfplots / CircuiTikz) sean consistentes entre clases.
%
% Es una COPIA de courses/Fisica3-2015/Clases/assets/tikzstyles.tex: mismo
% dominio, misma paleta. Copia y no symlink para que cada curso sea
% autocontenido y la paleta de Física III pueda evolucionar aparte.
%
% Uso: la clase debe cargar antes `tikz` y `pgfplots` (y `circuitikz` si la
%      clase tiene circuitos), y luego % courses/ElecMag2024/Clases/assets/tikzstyles.tex
% ---------------------------------------------------------------------------
% Estilos compartidos para los diagramas del curso Electromagnetismo (2024).
% Fuente única del "look" visual (colores, grosores, escalas) para que todas
% las figuras (TikZ / pgfplots / CircuiTikz) sean consistentes entre clases.
%
% Es una COPIA de courses/Fisica3-2015/Clases/assets/tikzstyles.tex: mismo
% dominio, misma paleta. Copia y no symlink para que cada curso sea
% autocontenido y la paleta de Física III pueda evolucionar aparte.
%
% Uso: la clase debe cargar antes `tikz` y `pgfplots` (y `circuitikz` si la
%      clase tiene circuitos), y luego \input{../assets/tikzstyles.tex}
% (compilar desde el directorio de la clase para que la ruta relativa resuelva).
% ---------------------------------------------------------------------------

% ---- Paleta (alineada con las cajas del preámbulo: keybox/notebox/warnbox) ----
\definecolor{figblue}{HTML}{1F4E79}   % azul principal (líneas, keybox)
\definecolor{figred}{HTML}{C0392B}    % rojo de énfasis (warnbox)
\definecolor{figamber}{HTML}{B8860B}  % ámbar (notebox)
\definecolor{figgray}{HTML}{555555}   % auxiliares, ejes, guías

% ---- CircuiTikz: escala y grosor de los componentes ----
% Guarda \ifdefined: la electrostática (Clases 1-14) no carga `circuitikz`, y
% sin la guarda el \input revienta con `Undefined control sequence \ctikzset`.
% Cuando el curso llegue a circuitos, el bloque se activa solo.
\ifdefined\ctikzset
  \ctikzset{
    bipoles/thickness=1,
    resistors/scale=0.9,
    capacitors/scale=0.9,
  }
\fi

% ---- TikZ: estilos reutilizables ----
% Nota: las puntas se declaran como `-{Latex[...]}` y no como `->`. La forma
% con llaves NO contiene un `>` literal, así que es inmune al `>` activo que
% introduce babel español (ver CLAUDE.md §7.3.3).
\usetikzlibrary{arrows.meta,calc,angles,quotes}

\tikzset{
  figline/.style={figblue, line width=0.9pt},
  figaux/.style={figgray, dashed, line width=0.6pt},
  % vectores
  figvec/.style={figblue, line width=0.9pt, -{Latex[length=2.2mm,width=1.6mm]}},
  figvecres/.style={figred, line width=1.3pt, -{Latex[length=2.6mm,width=1.9mm]}},
  figvecaux/.style={figgray, line width=0.7pt, -{Latex[length=1.8mm,width=1.3mm]}},
  figaxis/.style={figgray, line width=0.6pt, -{Latex[length=2mm,width=1.4mm]}},
  % cargas puntuales
  figcarga/.style={circle, inner sep=0pt, minimum size=5.4pt, draw=none},
  figpos/.style={figcarga, fill=figred},
  figneg/.style={figcarga, fill=figblue},
  figneutra/.style={figcarga, fill=figgray},
  % etiquetas
  figlbl/.style={font=\small},
  figlblaux/.style={font=\footnotesize, figgray},
}

% ---- pgfplots: estilo base de las gráficas ----
\pgfplotsset{
  emfig/.style={
    width=11cm, height=6.5cm,
    axis lines=left,
    xlabel style={font=\small},
    ylabel style={font=\small},
    tick label style={font=\footnotesize},
    every axis plot/.append style={line width=1pt},
    grid=major,
    grid style={figgray!20},
    legend cell align=left,
    legend style={font=\footnotesize, draw=figgray!40},
  },
}

% (compilar desde el directorio de la clase para que la ruta relativa resuelva).
% ---------------------------------------------------------------------------

% ---- Paleta (alineada con las cajas del preámbulo: keybox/notebox/warnbox) ----
\definecolor{figblue}{HTML}{1F4E79}   % azul principal (líneas, keybox)
\definecolor{figred}{HTML}{C0392B}    % rojo de énfasis (warnbox)
\definecolor{figamber}{HTML}{B8860B}  % ámbar (notebox)
\definecolor{figgray}{HTML}{555555}   % auxiliares, ejes, guías

% ---- CircuiTikz: escala y grosor de los componentes ----
% Guarda \ifdefined: la electrostática (Clases 1-14) no carga `circuitikz`, y
% sin la guarda el \input revienta con `Undefined control sequence \ctikzset`.
% Cuando el curso llegue a circuitos, el bloque se activa solo.
\ifdefined\ctikzset
  \ctikzset{
    bipoles/thickness=1,
    resistors/scale=0.9,
    capacitors/scale=0.9,
  }
\fi

% ---- TikZ: estilos reutilizables ----
% Nota: las puntas se declaran como `-{Latex[...]}` y no como `->`. La forma
% con llaves NO contiene un `>` literal, así que es inmune al `>` activo que
% introduce babel español (ver CLAUDE.md §7.3.3).
\usetikzlibrary{arrows.meta,calc,angles,quotes}

\tikzset{
  figline/.style={figblue, line width=0.9pt},
  figaux/.style={figgray, dashed, line width=0.6pt},
  % vectores
  figvec/.style={figblue, line width=0.9pt, -{Latex[length=2.2mm,width=1.6mm]}},
  figvecres/.style={figred, line width=1.3pt, -{Latex[length=2.6mm,width=1.9mm]}},
  figvecaux/.style={figgray, line width=0.7pt, -{Latex[length=1.8mm,width=1.3mm]}},
  figaxis/.style={figgray, line width=0.6pt, -{Latex[length=2mm,width=1.4mm]}},
  % cargas puntuales
  figcarga/.style={circle, inner sep=0pt, minimum size=5.4pt, draw=none},
  figpos/.style={figcarga, fill=figred},
  figneg/.style={figcarga, fill=figblue},
  figneutra/.style={figcarga, fill=figgray},
  % etiquetas
  figlbl/.style={font=\small},
  figlblaux/.style={font=\footnotesize, figgray},
}

% ---- pgfplots: estilo base de las gráficas ----
\pgfplotsset{
  emfig/.style={
    width=11cm, height=6.5cm,
    axis lines=left,
    xlabel style={font=\small},
    ylabel style={font=\small},
    tick label style={font=\footnotesize},
    every axis plot/.append style={line width=1pt},
    grid=major,
    grid style={figgray!20},
    legend cell align=left,
    legend style={font=\footnotesize, draw=figgray!40},
  },
}

% (compilar desde el directorio de la clase para que la ruta relativa resuelva).
% ---------------------------------------------------------------------------

% ---- Paleta (alineada con las cajas del preámbulo: keybox/notebox/warnbox) ----
\definecolor{figblue}{HTML}{1F4E79}   % azul principal (líneas, keybox)
\definecolor{figred}{HTML}{C0392B}    % rojo de énfasis (warnbox)
\definecolor{figamber}{HTML}{B8860B}  % ámbar (notebox)
\definecolor{figgray}{HTML}{555555}   % auxiliares, ejes, guías

% ---- CircuiTikz: escala y grosor de los componentes ----
% Guarda \ifdefined: la electrostática (Clases 1-14) no carga `circuitikz`, y
% sin la guarda el \input revienta con `Undefined control sequence \ctikzset`.
% Cuando el curso llegue a circuitos, el bloque se activa solo.
\ifdefined\ctikzset
  \ctikzset{
    bipoles/thickness=1,
    resistors/scale=0.9,
    capacitors/scale=0.9,
  }
\fi

% ---- TikZ: estilos reutilizables ----
% Nota: las puntas se declaran como `-{Latex[...]}` y no como `->`. La forma
% con llaves NO contiene un `>` literal, así que es inmune al `>` activo que
% introduce babel español (ver CLAUDE.md §7.3.3).
\usetikzlibrary{arrows.meta,calc,angles,quotes}

\tikzset{
  figline/.style={figblue, line width=0.9pt},
  figaux/.style={figgray, dashed, line width=0.6pt},
  % vectores
  figvec/.style={figblue, line width=0.9pt, -{Latex[length=2.2mm,width=1.6mm]}},
  figvecres/.style={figred, line width=1.3pt, -{Latex[length=2.6mm,width=1.9mm]}},
  figvecaux/.style={figgray, line width=0.7pt, -{Latex[length=1.8mm,width=1.3mm]}},
  figaxis/.style={figgray, line width=0.6pt, -{Latex[length=2mm,width=1.4mm]}},
  % cargas puntuales
  figcarga/.style={circle, inner sep=0pt, minimum size=5.4pt, draw=none},
  figpos/.style={figcarga, fill=figred},
  figneg/.style={figcarga, fill=figblue},
  figneutra/.style={figcarga, fill=figgray},
  % etiquetas
  figlbl/.style={font=\small},
  figlblaux/.style={font=\footnotesize, figgray},
}

% ---- pgfplots: estilo base de las gráficas ----
\pgfplotsset{
  emfig/.style={
    width=11cm, height=6.5cm,
    axis lines=left,
    xlabel style={font=\small},
    ylabel style={font=\small},
    tick label style={font=\footnotesize},
    every axis plot/.append style={line width=1pt},
    grid=major,
    grid style={figgray!20},
    legend cell align=left,
    legend style={font=\footnotesize, draw=figgray!40},
  },
}


% ---- Title ----
\title{\textbf{Clase 16: Campo magnético} \\[0.3em]
        \large{Fuerza entre cargas en movimiento, fuerza de Lorentz, fuerza y torque sobre circuitos}}
\author{Transcripto y expandido --- Curso Electromagnetismo (OpenFING)}
\date{}

\begin{document}

\maketitle
\thispagestyle{empty}
\newpage

\tableofcontents
\newpage

\section{Introducción histórica al magnetismo}

\subsection{De las magnetitas a las corrientes}

Que las cargas en movimiento interactúan con fuerzas adicionales a las
electrostáticas no es sorpresa; pero no fue así como se descubrió el
magnetismo, que se conocía mucho antes de saber qué eran las cargas eléctricas.

\begin{itemize}
  \item \textbf{Antigüedad.} Las \textbf{magnetitas} ---básicamente óxidos de hierro---, colgadas
  de un hilo o flotando en una cubeta, se orientan siempre en una misma
  dirección. Utilidad tecnológica: las brújulas. Interpretación durante buena
  parte de la historia: piedras mágicas.
  \item \textbf{Siglo XIX.} El primero en advertir la naturaleza oculta del fenómeno
  observó que el efecto de una \textbf{corriente eléctrica} es muy similar al de un
  \textbf{imán} (una piedra magnetizada).
\end{itemize}

\begin{notebox}
El VTT transcribe el nombre como «Hertz»/«Ertz»; el docente lo presenta como
«muy difícil de escribir» y lo escribe en el pizarrón, que no se ve. Por el
hecho histórico descrito se trata de \textbf{Oersted (Ørsted)}; el nombre no se
puede confirmar desde la transcripción.
\end{notebox}

De ahí salió una línea de investigación que ocupó todo el siglo XIX, con los
nombres que aparecerán en este capítulo y los siguientes: Ampère, Biot y
Savart, Faraday, Henry, Maxwell.

La conclusión, que es correcta: \textbf{hay un único magnetismo}. En los imanes hay,
en algún sentido, corrientes eléctricas «escondidas». El magnetismo es una
interacción entre \textbf{cargas en movimiento}; lo único que cambia es que en un
caso la corriente se ve fluir macroscópicamente y en el otro es interna.

\subsection{Magnetización y analogía con la polarización}

\begin{center}
\renewcommand{\arraystretch}{1.35}
\small
\begin{tabularx}{\textwidth}{>{\raggedright\arraybackslash}X>{\raggedright\arraybackslash}X}
\toprule
\textbf{Electricidad} & \textbf{Magnetismo} \\
\midrule
Cargas libres, a las que se tiene acceso & Corrientes eléctricas macroscópicas \\
Polarización & \textbf{Magnetización}: pequeñas corrientes a nivel atómico o molecular \\
\bottomrule
\end{tabularx}
\end{center}

Con una diferencia fundamental: la mayoría de los materiales \textbf{no} tiene
polarización espontánea (sin campo eléctrico externo, la polarización es cero),
mientras que muchos materiales \textbf{sí tienen magnetización espontánea}: se
mantienen como imanes sin que se les aplique ningún campo magnético.

\begin{notebox}
El docente comenta que una de sus dos líneas de investigación son justamente
las propiedades de materiales magnéticos, y que en el curso casi no se va a
hablar de eso.
\end{notebox}

\subsection{El espín: lo que no es clásico}

Una salvedad que, según el docente, ni siquiera menciona el Reitz: \textbf{no todas
esas corrientes son interpretables como cargas clásicas en movimiento}. En los
materiales magnéticos el efecto principal no viene de cargas que se desplazan
sino del \textbf{espín} de los electrones.

En una primera aproximación «con muchos defectos», el electrón sería una bolita
cargada que gira sobre sí misma: sus partes se mueven y hay una corriente. Pero
esa imagen no admite interpretación clásica, por una razón directa: la carga
está cuantizada y no hay carga más chica que la del electrón, así que no se lo
puede pensar como una bola con una distribución de carga adentro. Cuando se
intenta, se llega enseguida a absurdos.

\begin{notebox}[A nivel clásico no hay magnetismo de la materia]
La magnetización
espontánea tiene una componente cuántica fundamental; el docente afirma que
esto «se puede probar» y no lo prueba. Para el curso no importa el origen:
hay una magnetización que genera efectos magnéticos, y se tratará más
adelante.
\end{notebox}

Por ahora: interacción magnética entre cargas en movimiento y corrientes.

\section{Fuerza magnética entre dos cargas}

\subsection{La ley experimental}

En electrostática, dos cargas $q_1$ y $q_2$ interactúan según la ley de
Coulomb. Con $\vec r$ el vector que va de la carga 1 a la 2, la fuerza sobre
la 2 es

\[
\vec F_{e,2} = \frac{1}{4\pi\varepsilon_0}\,\frac{q_1 q_2}{r^3}\,\vec r .
\]

\begin{figure}[H]
  \centering
  % Clase 16 — dos cargas en movimiento. r va de la carga 1 a la 2; la fuerza
% magnética sobre la 2 hace intervenir las dos velocidades por separado (no la
% velocidad relativa), medidas en un referencial dado.
\begin{tikzpicture}[scale=1]

  \coordinate (q1) at (0,0);
  \coordinate (q2) at (4.6,1.6);

  \draw[figaux, -{Latex[length=2mm,width=1.4mm]}] (q1) -- (q2);
  \node[figlbl] at (2.2,1.12) {$\vec r$};

  \draw[figvec] (q1) -- (1.5,-0.75);
  \node[figlbl] at (1.8,-0.9) {$\vec v_1$};
  \draw[figvec] (q2) -- (5.5,2.85);
  \node[figlbl] at (5.75,3.05) {$\vec v_2$};

  \node[figpos] at (q1) {};
  \node[figpos] at (q2) {};
  \node[figlbl] at (-0.35,0.3) {$q_1$};
  \node[figlbl] at (4.95,1.25) {$q_2$};

\end{tikzpicture}

  \caption{Dos cargas en movimiento. $\vec r$ va de la carga 1 a la 2; las velocidades son las de cada carga en el referencial elegido.}
\end{figure}

Si las cargas se mueven, esa fuerza \textbf{sigue estando}, y se le agrega otra, la
\textbf{fuerza magnética}:

\begin{keybox}
\[
\boxed{\ \vec F_{m,2} = \frac{\mu_0}{4\pi}\,\frac{q_1 q_2}{r^3}\;
\vec v_2\times\left(\vec v_1\times\vec r\right)\ }
\]
\end{keybox}

Es una \textbf{ley experimental}. No es una alternativa a Coulomb sino un
complemento: la suma de las dos vale también en el caso estático, donde las
velocidades son nulas y el segundo término se anula. Hace falta que \textbf{las dos}
cargas se muevan.

\subsection{Dominio de validez y el problema del referencial}

\begin{warnbox}
\textbf{Esta forma sólo vale si $|\vec v_1|, |\vec v_2| \ll c$}, con
$c = 2{,}9979\times10^{8}$ m/s la velocidad de la luz. Las correcciones están
suprimidas por factores $v^2/c^2$: son órdenes superiores de un desarrollo en
$v/c$.
\end{warnbox}

Un estudiante pregunta lo obvio: \textbf{¿velocidades en qué referencial?} Y es «muy
raro». La mecánica newtoniana cumple el principio de Galileo ---en todo
referencial inercial valen las mismas leyes---, y esta fórmula hace referencia a
las velocidades en un referencial dado; peor, a las \textbf{dos velocidades por
separado} y no a la velocidad relativa. La ley \textbf{no es invariante} ante un
cambio de referencial.

Lo que el docente dice al respecto, sin desarrollarlo:

\begin{itemize}
  \item La fórmula vale en un referencial en que ambas partículas tengan velocidad
  chica; si no existe, no vale.
  \item Cuando se agregan los términos que la aproximación deja afuera, «todo se
  arregla». Entender cómo llevó un siglo: lo hizo Einstein en 1905, y es la
  teoría de la relatividad. No alcanzaba con corregir la fuerza
  electromagnética: hubo que cambiar la propia ley de Newton.
  \item «Por ahora no les estoy mostrando la solución, les estoy mostrando el
  problema.»
\end{itemize}

\begin{notebox}
Digresión histórica del docente: el principio lo formuló antes Giordano Bruno
(en la bodega de un gran barco en mar calmo, jugando a los dados y viendo
volar mariposas, no se puede saber si el barco se mueve); Galileo lo hizo
famoso.
\end{notebox}

\subsection{La constante \texorpdfstring{$\mu_0$}{mu0} y la velocidad de la luz}

$\mu_0$ juega para la fuerza magnética un papel análogo al de $\varepsilon_0$
para la eléctrica: la \textbf{permeabilidad magnética del vacío} (la transcripción
dice «permitividad magnética»). En el Sistema Internacional vale

\[
\frac{\mu_0}{4\pi} = 10^{-7}\ \frac{\text{N}\,\text{s}^2}{\text{C}^2}
\qquad\text{(exactamente)},
\]
sin ningún coeficiente raro: el sistema está elegido así, forma parte de la
definición del coulomb. Y las dos constantes están relacionadas:

\begin{keybox}
\[
\boxed{\ \mu_0\,\varepsilon_0 = \frac{1}{c^2}\ }
\]
\end{keybox}

Quien hizo Física III sabe que no es coincidencia (la luz es un fenómeno
electromagnético), pero por ahora se lo toma como un hecho útil.

\section{Campo magnético y fuerza de Lorentz}

Igual que en el caso eléctrico, en lugar de una acción a distancia conviene
pensar la fuerza como el efecto de un \textbf{campo} sobre una carga. En
electrostática se factoriza la carga que siente la fuerza:

\[
\vec F_{e,2} = q_2\,\vec E, \qquad
\vec E = \frac{q_1}{4\pi\varepsilon_0}\,\frac{\vec r}{r^3} .
\]

En el caso magnético hay \textbf{dos} propiedades de la carga 2 en juego, su carga y
su velocidad, y se factorizan las dos:

\[
\vec F_{m,2} = q_2\,\vec v_2\times\vec B,
\qquad
\boxed{\ \vec B = \frac{\mu_0}{4\pi}\,q_1\,\frac{\vec v_1\times\vec r}{r^3}\ }
\]
con $\vec B$ el \textbf{campo magnético} producido por la carga 1. La utilidad es la
misma que en el caso eléctrico: todo lo hecho con la ecuación de Laplace habría
sido muy difícil con interacción directa carga a carga.

Con campos eléctrico y magnético presentes, una carga siente la \textbf{fuerza de
Lorentz}:

\begin{keybox}
\[
\boxed{\ \vec F = q\left(\vec E + \vec v\times\vec B\right)\ }
\]
\end{keybox}

\begin{notebox}[¿Por qué un producto vectorial?]
«Porque no hay remedio, porque la fuerza
es así.» Y $\vec v_1$ en $\vec B$ \textbf{no} es una velocidad relativa: es la
velocidad de la carga 1 en el referencial elegido.
\end{notebox}

\section{Comparación entre la fuerza magnética y la eléctrica}

\subsection{La cota}

Usando $\mu_0 = 1/(\varepsilon_0 c^2)$ y repartiendo un $c$ a cada velocidad:

\[
\vec F_{m,2} = \frac{1}{4\pi\varepsilon_0}\,\frac{q_1 q_2}{r^3}\;
\frac{\vec v_2}{c}\times\left(\frac{\vec v_1}{c}\times\vec r\right).
\]

Al tomar módulos aparecen senos, que son menores o iguales que 1:

\begin{keybox}
\[
\boxed{\ |\vec F_{m,2}| \;\le\; |\vec F_{e,2}|\;\frac{v_1}{c}\,\frac{v_2}{c}\ }
\]
\end{keybox}

Para los fenómenos cotidianos ---un auto a 300 km/h es insignificante frente a
$c$--- la fuerza magnética entre dos cargas es \textbf{muchísimo} menor que la
eléctrica. El factor va entre 0 y 1 y, a velocidades normales, está «pegado a
cero».

\subsection{Por qué igual se observa el magnetismo}

Si fuera sólo eso, los fenómenos magnéticos serían inobservables. No lo son
porque no se ven en pares de cargas sino en \textbf{conjuntos grandes de cargas con
carga neta nula y corriente no nula}.

\begin{notebox}[Ejemplo clave: un cable con corriente]
El cable es neutro, no ejerce
fuerza de Coulomb; pero la corriente no es cero y el efecto magnético
sobrevive. Para ver un efecto magnético a velocidades normales hace falta que
el efecto electrostático sea cero.
\end{notebox}

Recordatorio a raíz de una pregunta: la velocidad individual de cada electrón
en un conductor es muy alta, pero la velocidad de deriva es chica («algunos
centímetros por hora», según el material).

\section{Fuerza sobre un cable con corriente}

\subsection{Un tramo de cable}

Hay un campo magnético $\vec B$, producido por algo que por ahora no interesa, y
un cable descrito por una curva $C$ con corriente $I$. Se piensa el cable como
un cañito de sección $A$. En un tramito $d\vec l$ la velocidad de deriva es
\textbf{paralela} a $d\vec l$.

\begin{figure}[H]
  \centering
  % Clase 16 — fuerza sobre un tramo de cable. El cable es un cañito de sección
% A; en el tramo dl la velocidad de deriva es paralela a dl. Con dl hacia la
% derecha y B hacia arriba en el plano del papel, dF = I dl x B sale del papel.
\begin{tikzpicture}[scale=1]

  % campo B de fondo
  \foreach \bx in {1.9,5.3,6.5} {
    \draw[figvecaux] (\bx,0.95) -- (\bx,1.85);
  }
  \foreach \bx in {0.7,1.9,6.9} {
    \draw[figvecaux] (\bx,-1.35) -- (\bx,-0.45);
  }
  \node[figlblaux] at (6.85,1.5) {$\vec B$};

  % cable
  \fill[figgray!14] (0,-0.4) rectangle (7.2,0.4);
  \draw[figline] (0,0.4) -- (7.2,0.4);
  \draw[figline] (0,-0.4) -- (7.2,-0.4);
  \draw[figline, fill=figgray!28] (7.2,0) ellipse (0.16 and 0.4);
  \draw[figline] (0,0.4) arc (90:270:0.16 and 0.4);
  \node[figlbl] at (7.65,0) {$A$};

  % tramo dl
  \fill[figblue!22] (2.9,-0.4) rectangle (4.3,0.4);
  \draw[figaux] (2.9,-0.4) -- (2.9,0.4);
  \draw[figaux] (4.3,-0.4) -- (4.3,0.4);
  \foreach \px/\py in {3.05/0.17, 3.2/-0.18} {
    \node[figneg] at (\px,\py) {};
    \draw[figvec] (\px,\py) -- ++(0.7,0);
  }
  \node[figlblaux] at (4.75,-0.17) {$\vec v_d$};

  % dl y corriente
  \draw[figvecres] (2.9,0.72) -- (4.3,0.72);
  \node[figlbl, text=figred] at (3.6,1.05) {$d\vec l$};
  \draw[figvecres] (0.2,0.72) -- (1.1,0.72);
  \node[figlbl, text=figred] at (0.65,1.05) {$I$};

  % dF saliente
  \draw[figred, line width=0.9pt] (3.6,-0.95) circle (0.17);
  \fill[figred] (3.6,-0.95) circle (0.05);
  \node[figlbl, text=figred, anchor=west] at (3.9,-0.95) {$d\vec F = I\,d\vec l\times\vec B$};

\end{tikzpicture}

  \caption{Tramo $d\vec l$ de un cable de sección $A$ en un campo magnético. Con $d\vec l$ hacia la derecha y $\vec B$ hacia arriba, la fuerza sobre el tramo sale del papel.}
\end{figure}

La fuerza sobre el tramito es la fuerza sobre un portador por el número de
portadores que hay en él. Con $n$ portadores por unidad de volumen (un solo
tipo de portador, como en un metal; con varios se suma sobre los tipos):

\[
d\vec F = \underbrace{q\,\vec v_d\times\vec B}_{\text{un portador}}\;
\underbrace{n\,A\,dl}_{\text{portadores en } dl}
= A\,n\,q\,dl\;\vec v_d\times\vec B .
\]

Como $\vec v_d$ y $d\vec l$ tienen igual dirección y sentido,
$dl\,\vec v_d = v_d\,d\vec l$: se pasa el carácter vectorial de uno al otro.
Y $n\,q\,v_d$ es la densidad de corriente, que por el área da la corriente
($I = JA$ en un conductor filiforme):

\begin{keybox}
\[
\boxed{\ d\vec F = I\,d\vec l\times\vec B\ }
\]
\end{keybox}

\begin{notebox}
Si además hubiera campo eléctrico, el cable no sentiría fuerza eléctrica: es
neutro. Sólo sobrevive la magnética.
\end{notebox}

Para un circuito completo se suma sobre la curva; la corriente es la misma en
todo el circuito y sale de la integral:

\[
\vec F = I\oint_C d\vec l\times\vec B .
\]

\subsection{Espira cerrada en un campo uniforme}

Si $\vec B$ es uniforme sale de la integral ---\textbf{por la derecha}, porque el
producto vectorial no es conmutativo---:

\[
\vec F = I\left(\oint_C d\vec l\right)\times\vec B = 0,
\]
porque la suma de todos los $d\vec l$ a lo largo de una curva cerrada es nula.

\begin{keybox}
\[
\boxed{\ \text{Una espira en un campo magnético uniforme no experimenta fuerza neta.}\ }
\]
\end{keybox}

Para que haya fuerza el campo tiene que ser \textbf{no uniforme}, de modo que partes
distintas de la espira sientan fuerzas distintas. Es la situación normal: cerca
de un imán el campo es más intenso que lejos, y los lados de la espira más
próximos no se compensan con los otros.

\section{Torque sobre una espira en un campo uniforme}

Aun en un campo uniforme, una espira \textbf{sí} experimenta un torque. En Física III
se vio para una espira plana rectangular; ahora se hace para \textbf{cualquier curva
cerrada}, plana o no. (Si el campo no es uniforme hay las dos cosas, fuerza y
torque.)

\subsection{Planteo en componentes}

Respecto de un origen, el torque de un elemento de fuerza es
$d\vec\tau = \vec r\times d\vec F$:

\[
\vec\tau = I\oint_C \vec r\times\left(d\vec l\times\vec B\right).
\]

Con $d\vec l = (dx, dy, dz)$ y $\vec B = (B_x, B_y, B_z)$, por la regla del
determinante:

\[
d\vec l\times\vec B =
(dy\,B_z - dz\,B_y)\,\hat\imath
+ (dz\,B_x - dx\,B_z)\,\hat\jmath
+ (dx\,B_y - dy\,B_x)\,\hat k .
\]

Se calcula sólo la componente $x$; las otras salen por \textbf{permutación circular}.
Con $\vec r = (x,y,z)$:

\[
\left[\vec r\times(d\vec l\times\vec B)\right]_x
= y\,(dx\,B_y - dy\,B_x) - z\,(dz\,B_x - dx\,B_z).
\]

\subsection{Las cuatro integrales}

Como $\vec B$ es uniforme, sus componentes salen de las integrales:

\[
\frac{\tau_x}{I} = B_y\oint_C y\,dx \;-\; B_x\oint_C y\,dy
\;-\; B_x\oint_C z\,dz \;+\; B_z\oint_C z\,dx .
\]

\textbf{Las dos fáciles.} $\oint y\,dy = 0$ y $\oint z\,dz = 0$: en una curva cerrada
se integra de $y_{\min}$ a $y_{\max}$ por un lado y se vuelve por el otro; lo
que se suma a la ida se deshace a la vuelta.

\textbf{Las otras dos.} $\oint_C y\,dx$ no hace intervenir $z$, así que es idéntica
a la integral sobre $C'$, la \textbf{proyección} de $C$ en el plano $Oxy$. Ya es un
problema plano. Se describe $C'$ con dos funciones, $y_1(x)$ por abajo e
$y_2(x)$ por arriba, entre $x_{\min}$ y $x_{\max}$:

\[
\oint_{C'} y\,dx = \int_{x_{\min}}^{x_{\max}} y_1\,dx - \int_{x_{\min}}^{x_{\max}} y_2\,dx
= -\,\text{área}(C') .
\]

\begin{notebox}
No vale cero porque el integrando no tiene función potencial: $y$ cambia a
lo largo de la curva.
\end{notebox}

\begin{figure}[H]
  \centering
  % Clase 16 — la integral de y dx sobre la curva proyectada C'. Se recorre por
% abajo (y_1) de x_min a x_max y se vuelve por arriba (y_2): al restar queda
% menos el área encerrada.
\begin{tikzpicture}[scale=1]

  % region
  \fill[figblue!12]
    plot[smooth cycle, tension=0.75] coordinates {(1,1.75) (2,0.8) (3.6,0.65) (4.9,1.35) (5,2.3) (4,3.05) (2.5,3.1) (1.3,2.6)};
  \draw[figline, line width=1pt]
    plot[smooth cycle, tension=0.75] coordinates {(1,1.75) (2,0.8) (3.6,0.65) (4.9,1.35) (5,2.3) (4,3.05) (2.5,3.1) (1.3,2.6)};

  % ejes
  \draw[figaxis] (-0.2,0) -- (6.2,0);
  \draw[figaxis] (0,-0.2) -- (0,3.9);
  \node[figlbl] at (6.45,0) {$x$};
  \node[figlbl] at (0,4.15) {$y$};

  % extremos
  \draw[figaux] (0.93,0) -- (0.93,2.0);
  \draw[figaux] (5.1,0) -- (5.1,1.9);
  \node[figlblaux] at (0.93,-0.32) {$x_{\min}$};
  \node[figlblaux] at (5.1,-0.32) {$x_{\max}$};

  % sentido de recorrido
  \draw[figvecres] (2.6,0.66) -- (3.3,0.62);
  \draw[figvecres] (3.4,3.13) -- (2.7,3.12);

  % rotulos
  \node[figlbl] at (3.0,0.28) {$y_1(x)$};
  \node[figlbl] at (3.05,3.5) {$y_2(x)$};
  \node[figlbl] at (3.0,1.9) {área de $C'$};
  \node[figlbl] at (5.55,2.75) {$C'$};

\end{tikzpicture}

  \caption{La curva proyectada $C'$ en el plano $Oxy$. Recorrida por abajo, $y_1(x)$, y de vuelta por arriba, $y_2(x)$, la integral de $y\,dx$ da menos el área encerrada.}
\end{figure}

\subsection{El vector área y el resultado}

Se definen las tres áreas de las proyecciones de $C$:

\begin{center}
\renewcommand{\arraystretch}{1.35}
\small
\begin{tabularx}{\textwidth}{l>{\raggedright\arraybackslash}X>{\raggedright\arraybackslash}X}
\toprule
\textbf{Componente} & \textbf{Proyección de $C$ en el plano} & \textbf{Integral} \\
\midrule
$A_z$ & $Oxy$ & $\oint x\,dy = -\oint y\,dx$ \\
$A_x$ & $Oyz$ & $\oint y\,dz = -\oint z\,dy$ \\
$A_y$ & $Ozx$ & $\oint z\,dx = -\oint x\,dz$ \\
\bottomrule
\end{tabularx}
\end{center}

Lo que se mantiene es la permutación circular $x\to y\to z\to x$; invertir el
orden de las dos coordenadas cambia el signo, porque la proyección se recorre
en el sentido opuesto. Con eso, $\oint y\,dx = -A_z$ y $\oint z\,dx = A_y$:

\[
\tau_x = I\left(A_y B_z - A_z B_y\right) = I\left(\vec A\times\vec B\right)_x .
\]

Y lo mismo para las otras componentes:

\begin{keybox}
\[
\boxed{\ \vec\tau = I\,\vec A\times\vec B\ }
\qquad\text{(curva cerrada cualquiera, } \vec B \text{ uniforme)}
\]
\end{keybox}
con $\vec A = (A_x, A_y, A_z)$ el \textbf{vector área}.

\begin{notebox}[Qué hay que retener]
El resultado es el mismo de Física III, pero vale
para curvas en el espacio. En el caso plano el vector área era «el área por
la normal»; para una curva no plana no es obvio qué significa, y la respuesta
es: el vector cuyas componentes cartesianas son las áreas de las proyecciones
en los tres planos.
\end{notebox}

\begin{notebox}
El docente tropieza varias veces con los signos en el pizarrón y deja la
verificación como tarea. A la pregunta de cómo se orienta el área responde
que las componentes son áreas «en el sentido usual» y que el vector queda sin
ambigüedad; el sentido lo fijan las integrales de la tabla, que dependen del
sentido en que se recorre la curva (el de la corriente).
\end{notebox}

\textbf{Ejercicio propuesto} (generalizar el razonamiento componente a componente):

\[
\vec A = \frac12\oint_C \vec r\times d\vec l .
\]

\begin{notebox}
La transcripción dice «$\vec r$ escalar $d\vec l$»; para que el resultado sea
un vector, y coincida con la tabla, tiene que ser el producto vectorial.
\end{notebox}

\subsection{Momento dipolar magnético}

A la combinación

\begin{keybox}
\[
\boxed{\ \vec m = I\,\vec A\ }
\]
\end{keybox}
se la llama \textbf{momento dipolar magnético}, de modo que
$\vec\tau = \vec m\times\vec B$.

\section{Corriente distribuida en un volumen}

Todo lo anterior es para cables. Se generaliza a una densidad de corriente
$\vec J(\vec r)$ distribuida en un volumen dividiéndola en \textbf{cablecitos
filiformes imaginarios} que siguen las líneas de corriente. En un tramito de
uno de ellos, de sección $A$, se toma $d\vec l$ paralelo a $\vec J$, así que
$I\,d\vec l = A\,J\,d\vec l = A\,|d\vec l|\,\vec J$:

\[
d\vec F = I\,d\vec l\times\vec B = A\,|d\vec l|\;\vec J\times\vec B
= \vec J\times\vec B\;dV,
\]
porque $A\,|d\vec l|$ es el elemento de volumen. En lugar de integrar sobre una
curva se integra sobre el volumen:

\begin{keybox}
\[
\boxed{\ \vec F = \int_V \vec J\times\vec B\;dV\ }
\]
\end{keybox}

Regla para recordarla: donde había $I\,d\vec l$ va $\vec J\,dV$. Con el torque
se puede hacer lo mismo, sin sorpresas; el docente no lo escribe.

\noindent\emph{La próxima clase: la ley de Biot y Savart ---cómo una distribución de corrientes
genera un campo magnético--- y, lo que en Física III se afirmaba sin probar, la
equivalencia entre Biot--Savart y la ley de Ampère, para lo cual hace falta
cálculo vectorial.}

\end{document}
